\section{The relative ray-class rung and the cubic arithmetic gate}
\label{sec:kummer-lambda}

We use the newest supplied ALP\_MACHINE checkpoint, dated
2026-11-21T2304Z. Its AM-H01--AM-H04 identify the cubic representation,
its full Euler factors, the obstruction to a cyclotomic or CM overfield,
and the analytic window implication. They are inputs here, with the
checks recorded separately. In particular, we do not announce those
calculations as new findings. Put
\[
 K=\mathbb Q(\omega),\quad \omega^2+\omega+1=0,
 \quad E_p=K(\sqrt[3]p),\quad F_p=\mathbb Q(\sqrt[3]p),
 \quad L_p(s)=L_K(s,\chi_{p,3}).
\]
The character is understood primitively, with every local factor
defined on inertia invariants. AM-H01 gives
\begin{equation}
 L_p(s)=\zeta_{F_p}(s)/\zeta(s),\qquad
 \zeta_{E_p}(s)=\zeta_K(s)L_p(s)^2.                 \label{eq:kummer-zeta}
\end{equation}
These are identities of meromorphic functions, including the bad primes.

\subsection{The precise meaning of integral Lambda descent}

\paragraph{Known classification, with its actual scope.}
Borger--de Smit, \emph{Galois theory and integral models of
$\Lambda$-rings}, Theorem~1.1 (source label \texttt{thm-A}),
classifies finite \'etale \emph{algebras}, with their commuting
Frobenius lifts and an integral model, by finite continuous sets
for the multiplicative monoid $\widehat{\mathbb Z}$.
The paragraph immediately following that theorem explicitly says
that such an algebra cannot be a field other than $\mathbb Q$.
Theorem~1.2 (\texttt{thm-B}) identifies the maximal integral model
of $\mathbb Q[t]/(t^n-1)$ as $\mathbb Z[t]/(t^n-1)$.
For a number field base, Borger--de Smit, \emph{Lambda actions of
rings of integers}, Theorems~1.1 and~1.2 (\texttt{local-thm},
\texttt{global-thm}), replace the monoid by the Deligne--Ribet
monoid. The source has the typographical phrase ``Then $K$ has an
integral model'' in both theorem statements; their quantified
algebra and proofs make clear that $E$ is intended.

\begin{proposition}[AC-K01: the corrected rational equivalence]
\label{AC-K01}
The following zero-location assertions are equivalent:
\begin{enumerate}
\item input theorem A;
\item RH for every cyclotomic Dedekind zeta function;
\item RH for the product of Dedekind zeta functions of the field
factors of every finite \'etale $\mathbb Q$-algebra admitting an
integral $\Lambda_{\mathbb Z}$-model.
\end{enumerate}
Every abelian number field occurs as a field factor of an algebra
in (3). The assertion that the field itself has such a model is
false unless it is $\mathbb Q$.
\end{proposition}
\begin{proof}
Abelian Artin factorization expresses a cyclotomic Dedekind zeta
function as a product of primitive Dirichlet $L$-functions.
Every factor is holomorphic in $0<\Re s<1$, so an off-line zero
of one factor cannot be canceled by a pole of another. This proves
both (1)$\Rightarrow$(2) and (2)$\Rightarrow$(1): each primitive
Dirichlet character occurs at its conductor. Borger--de Smit's
classification makes each field factor in (3) abelian over
$\mathbb Q$, proving (1)$\Rightarrow$(3) by the same factorization.
For the converse use $\mathbb Q[t]/(t^n-1)$, whose factors include
$\mathbb Q(\mu_n)$; the same no-cancellation argument proves (2).

For the additional assertion, write the abelian field as
$F=\mathbb Q(\mu_n)^H$. The invariant subring
$\bigl(\mathbb Z[t]/(t^n-1)\bigr)^H$ is finite torsion free over
$\mathbb Z$ and is preserved by all maps $t\mapsto t^p$.
The Frobenius congruence remains valid in the invariant ring:
if an invariant element is $p$ times an element of the larger
torsion-free ring, that quotient is invariant too. Its generic
algebra has $F$ as the factor from primitive $n$th roots. The final
claim is exactly the cited paragraph after Theorem~1.1; the
relative elementary obstruction below also proves it.
\end{proof}

\begin{proposition}[AC-K02: a field obstruction at ramification]
\label{AC-K02}
Let $L/K$ be a finite field extension. If $L$ admits an integral
$\Lambda_{\mathcal O_K}$-model, it is unramified at every finite
prime of $K$. In particular no $E_p$ above admits such a model as
a field, even though $E_p/K$ is cyclic. This refutes the literal
field version of item~4 of the initial first-pass note.
\end{proposition}
\begin{proof}
Let $R\subset L$ be the finite flat integral model. Every
$K$-endomorphism $\psi_{\mathfrak p}$ of $L$ is an automorphism
of finite order. Since it preserves $R$, its inverse, a positive
power of itself, preserves $R$ too. Thus its reduction modulo
$\mathfrak p$ is an automorphism. By the defining congruence that
reduction is $x\mapsto x^{N\mathfrak p}$. A finite-dimensional
algebra over a finite field has injective Frobenius only if it
is reduced: a nonzero nilpotent would be killed by an iterate.
A reduced finite algebra over a perfect field is \'etale.
Hence $R$ is finite flat with \'etale geometric fibres, so it is
finite \'etale over $\mathcal O_K$. Consequently its fraction
field is unramified at every finite place.

For $p\ne3$, the valuation of $p$ at each prime of $K$ above $p$
is one; $X^3-p$ is Eisenstein there and $E_p/K$ is ramified.
For $p=3$, at $\mathfrak l=(1-\omega)$ the valuation of $3$ is
two, so the valuation of its cube root is $2/3$; the ramification
index is divisible by three. This also excludes that case.
\end{proof}

\paragraph{Geometric translation.}
AC-K01 concerns generic fibres of finite flat schemes over
$\operatorname{Spec}\mathbb Z$ with commuting Frobenius lifts.
The field factors need not be stable under these lifts.
AC-K02 is the distinction between a ramified cover and a finite
\'etale cover of $\operatorname{Spec}\mathcal O_K$, not an
obstruction to a rank-one lisse sheaf on an open subscheme.
The Galois-category framework is SGA~1, expos\'e~V, Theorem~4.1;
the arithmetic classification is the cited Borger--de Smit theorem,
not a theorem stated in SGA~1.

\subsection{An actual finite host, and its integral construction}

\begin{proposition}[AC-K03: the minimal-conductor relative host]
\label{AC-K03}
Let $\mathfrak f$ be the conductor of $\chi_{p,3}$ and let
$d(\mathfrak f)$ denote the number of integral-ideal divisors
of $\mathfrak f$. There is a finite \'etale $K$-algebra with
integral $\Lambda_{\mathcal O_K}$-model whose underlying algebra is
\begin{equation}
 B_{p}=E_p\times K^{d(\mathfrak f)-1}.              \label{eq:minhost}
\end{equation}
Its Frobenius lifts and an algorithm terminating at its maximal
integral model are constructed below. Its Dedekind zeta product is
\begin{equation}
 Z(B_p,s)=\zeta_K(s)^{d(\mathfrak f)}L_p(s)^2.       \label{eq:minhost-zeta}
\end{equation}
Under A, RH for $B_p$ is equivalent to RH for $L_p$.
Thus existence of this relative model does not remove the cubic
zero-location question.
\end{proposition}
\begin{proof}
Write $G=\operatorname{Cl}_{\mathfrak f}(K)$ and let
$H=\ker(G\twoheadrightarrow\operatorname{Gal}(E_p/K))$.
The finite set $S=\operatorname{DR}_{\mathfrak f}/H$ is a set
of orbits for multiplication by $H$. Commutativity makes
multiplication by every ideal a well-defined action on $S$.
It commutes with the Galois action through Artin reciprocity.
Set
\[
 B_p=\{f:S\longrightarrow\overline K:
                 f(gs)=g f(s)\text{ for }g\in G_K\},\qquad
 (\psi_{\mathfrak a}f)(s)=f(\mathfrak a s).
\]
This is an explicit finite \'etale algebra with commuting
$K$-algebra endomorphisms. The Deligne--Ribet factorization
criterion gives its integral model.

The decomposition
$\operatorname{DR}_{\mathfrak f}=
 \coprod_{\mathfrak d\mid\mathfrak f}
 \operatorname{Cl}_{\mathfrak f/\mathfrak d}(K)$
is stated immediately before Borger--de Smit's global theorem;
Yalkinoglu's equation \texttt{arithalg} gives the corresponding
product of ray class fields. After quotienting by $H$, the field
at $\mathfrak d$ is
$E_p\cap K_{\mathfrak f/\mathfrak d}$, by ordinary Galois
correspondence and the surjection of ray class groups.
For $\mathfrak d=1$ it is $E_p$. For $\mathfrak d\ne1$ it is
$K$: a degree-three extension has no proper intermediate field,
and containment in the smaller ray class field would contradict
minimality of its conductor. This proves \eqref{eq:minhost}.
Equation \eqref{eq:kummer-zeta} proves \eqref{eq:minhost-zeta}.
All factors are holomorphic in the open critical strip, so A and
the same no-cancellation argument prove the last assertion.

Here is a terminating integral construction, to avoid leaving
``has a model'' as merely an existential instruction. Put
$R_0=\mathcal O_{E_p}\times\mathcal O_K^{d(\mathfrak f)-1}$,
and let $T$ be the finite set of primes dividing $\mathfrak f$.
Recursively compute the subring
\begin{equation}
 R_{n+1}=\{x\in R_n:
 \psi_{\mathfrak p}(x)-x^{N\mathfrak p}
       \in\mathfrak p R_n\text{ for all }\mathfrak p\in T\}.
                                                               \label{eq:integral-iteration}
\end{equation}
It is an $\mathcal O_K$-subalgebra: modulo $\mathfrak p$ the
power map is additive, multiplicative and $\mathcal O_K/\mathfrak p$-linear.
All $\psi$ preserve $R_0$, because they preserve integrality.
Their commutativity proves by induction that they preserve every
$R_n$: apply $\psi_{\mathfrak q}$ to each congruence in
\eqref{eq:integral-iteration}. Each step is a finite computation
of kernels in $R_n/\mathfrak p R_n$, hence of a finite-index
submodule. Outside $T$ the original Frobenius congruences hold
in $R_0$, and each index $[R_0:R_n]$ is supported on $T$.
Consequently they also hold in $R_n$, since localization at
$\mathfrak q\notin T$ leaves $R_n$ unchanged.

By the classification there is some integral model $R$ of full
rank inside $R_0$. Its Frobenius congruences show inductively
that $R\subset R_n$ for every $n$. There are only finitely many
submodules between $R$ and $R_0$. The descending sequence
therefore reaches a fixed point. At that point all congruences
hold at all primes, so it is an integral model. The same argument
for any integral model shows it is the largest one. This proves
termination without asserting a uniform bound on its running
time or its index. Precisely such quantitative bounds would be
additional information in an arithmetic window argument.
\end{proof}

\begin{proposition}[AC-K04: the explicit host for $p=269$]
\label{AC-K04}
For $p=269$, the conductor of $\chi_{269,3}$ is the prime
$\mathfrak p=(269)$ of $\mathcal O_K$, of norm $269^2$.
Let $\mathfrak P$ be the unique prime of $E=E_{269}$ above
$\mathfrak p$. An explicit maximal integral model of
$B=E\times K$ is
\begin{equation}
 R=\{(x,y)\in\mathcal O_E\times\mathcal O_K:
                  x\bmod\mathfrak P=y\bmod\mathfrak p\}.
                                                               \label{eq:269-model}
\end{equation}
Its lifts are
\[
 \psi_{\mathfrak p}(x,y)=(y,y),\qquad
 \psi_{\mathfrak q}(x,y)=
       (\operatorname{Frob}_{\mathfrak q}x,y)
                     \quad(\mathfrak q\ne\mathfrak p),
\]
where arithmetic Frobenius is used. The index
$[\mathcal O_E\times\mathcal O_K:R]$ is $269^2$.
\end{proposition}
\begin{proof}
The prime is inert in $K$ since $269\equiv2\pmod3$.
The polynomial $X^3-269$ is Eisenstein at $\mathfrak p$, so
$E/K$ is totally and tamely ramified of degree three there.
All nontrivial characters of this tame inertia group have
conductor exponent one. The polynomial discriminant restricts
other possible ramification to the prime above three.

Since $269=-1+9\cdot30$, the equation
$t-3t^2+3t^3=30$ has a solution in $\mathbb Z_3$:
its derivative is one modulo three, and Hensel lifting applies.
Then $(-1+3t)^3=269$. As $\mu_3\subset K_3$, all three roots
are in $K_3$, so $E/K$ splits at the prime above three.
This proves the conductor assertion. The ray-class exact
sequence and class number one of $\mathbb Z[\omega]$ give
\[
 \operatorname{Cl}_{\mathfrak p}(K)
   \simeq\mathbb F_{269^2}^{\times}/\mu_6,
       \qquad |\operatorname{Cl}_{\mathfrak p}(K)|=12060.
\]
The group is cyclic and has one quotient of order three. Thus
the quotient Deligne--Ribet set in AC-K03 is a free $C_3$-orbit
and one fixed point. A prime away from $\mathfrak p$ acts on
the orbit by its Artin symbol; $\mathfrak p$ sends every point
to the fixed point. This gives exactly the displayed lifts.

The congruence condition in \eqref{eq:269-model} is stable under
these lifts, and the lifts commute. Away from $\mathfrak p$,
the algebra of integers is \'etale and the specified Frobenius
automorphism lifts residue Frobenius. Since the index of $R$
is supported at $\mathfrak p$, the same congruence holds in $R$.
At $\mathfrak p$ put $q=269^2$ and $x=y+t$ with
$t\in\mathfrak P$. We have $\mathfrak p\mathcal O_E=\mathfrak P^3$.
Every intermediate binomial coefficient of $(y+t)^q$ is
divisible by $269$, and $q\ge4$, so
\[
                 x^q-y^q\in\mathfrak P^4
                                      =\mathfrak p\mathfrak P.
\]
Also $y-y^q\in\mathfrak p$. Therefore
\[
 \frac{\psi_{\mathfrak p}(x,y)-(x,y)^q}{269}
 =\left(\frac{y-x^q}{269},\frac{y-y^q}{269}\right)\in R,
\]
because the difference of its coordinates belongs to
$\mathfrak P$. This is the required Frobenius congruence.
The residue-difference map from the maximal order onto
$\mathbb F_{269^2}$ has kernel $R$, proving its index.

Finally any integral model lies in the maximal order. Its
Frobenius congruence at $\mathfrak p$, reduced in the first
coordinate modulo $\mathfrak P$, forces $y=x^q=x$ in
$\mathbb F_{269^2}$. Hence it is contained in $R$, which proves
maximality. No bound on a Selmer group or on H04.4 has been
deduced from this exact index calculation.
\end{proof}

\subsection{Which Hecke sectors A actually supplies}

\begin{proposition}[AC-K05: norm sectors and infinity types]
\label{AC-K05}
For a primitive Dirichlet character $\eta$, the Hecke character
$\eta\circ N_{K/\mathbb Q}$ satisfies
\begin{equation}
 L_K(s,\eta\circ N)=L_{\mathbb Q}(s,\eta)
                    L_{\mathbb Q}(s,\eta\chi_{-3}),             \label{eq:norm-sectors}
\end{equation}
with primitive Artin local factors on both sides. A therefore
proves its critical-line assertion, and that of each unitary
norm twist. Every finite-order character of $G_K$ invariant
under complex conjugation is of this form.

A unitary Hecke character at the complex place has the form
$(z/|z|)^m|z|^{2it}$, with $m\in\mathbb Z$. Its finite component
on $\mu_6$ must be $\epsilon\mapsto\epsilon^{-m}$.
If the finite component is unramified everywhere, $6\mid m$;
ramification allows other $m$. The algebraic Hecke characters
of CM elliptic curves have nonzero angular type after unitary
normalization and are not norm sectors in \eqref{eq:norm-sectors}.
\end{proposition}
\begin{proof}
Induction of the restriction of the one-dimensional
$G_{\mathbb Q}$-representation $\eta$ is
$\eta\otimes(1\oplus\chi_{-3})$. Artin induction, including
inertia invariants, gives \eqref{eq:norm-sectors}. Tensoring by
the unitary norm character shifts $s$ vertically and leaves
the real critical line unchanged. If a finite-order character
$\xi$ of $G_K$ is conjugation invariant, the exact sequence
$G_K\subset G_{\mathbb Q}$ is split by an actual complex
conjugation $c$, with $c^2=1$. The formula
$\widetilde\xi(gc^e)=\xi(g)$, $e=0,1$, defines a continuous
finite-order character of $G_{\mathbb Q}$. By Kronecker--Weber
it is a Dirichlet character, and reciprocity identifies its
restriction with the norm pullback. Conversely a restriction
is conjugation invariant.

The continuous unitary characters of $\mathbb C^\times$
are exactly the displayed characters, by
$\mathbb C^\times\simeq\mathbb R_{>0}\times S^1$.
Triviality on principal global units gives the assertion on
$\mu_6$. A finite unramified component is trivial on all
finite local units, whence $6\mid m$. Finally complex norm
has no angular dependence, whereas a CM elliptic $H^1$
has types $(1,0)$ and $(0,1)$; its unitary characters have
angular exponents $1$ and $-1$ up to the reciprocity convention.
Thus they cannot be obtained by the specified norm pullback.
\end{proof}

\begin{proposition}[AC-K06: CM base change does not repair the slot]
\label{AC-K06}
Let $M$ be any CM number field containing $K$. The restriction
of $\chi_{p,3}$ to $G_M$ is nontrivial and cannot become the
restriction of a Dirichlet character of $G_{\mathbb Q}$.
Twisting by such a restriction does not change this obstruction.
\end{proposition}
\begin{proof}
AM-H03 proves that no CM field contains $E_p$.
Since $[E_p:K]=3$, its intersection with $M$ is $K$, so the
restricted character cuts out $E_pM/M$ and remains nontrivial.
If it were the restriction of a Dirichlet character of conductor
$n$, this field would lie in $M(\mu_n)$. A compositum of CM
fields is CM: in a common normal closure, complex conjugation
is central and restricts to the CM involution on each factor;
its fixed field is totally real and its nontrivial extension
has degree two. Thus $M(\mu_n)$ would be a CM field containing
$E_p$, contrary to AM-H03. A putative equality after a
Dirichlet twist can be divided by that twist and gives the
same contradiction.
\end{proof}

\subsection{The Bost--Connes realization and the missing positivity}

\paragraph{Known constructions.}
Laca--Larsen--Neshveyev, Theorem~2.1 (\texttt{mainthm}),
constructs and classifies the KMS states of the number-field
system, with the partition function $\zeta_K$ and the explicit
state formula \texttt{eKMSext}. Yalkinoglu, Theorems~A and~B
(\texttt{ThmA}, \texttt{ThmB}), identifies its measured
endomotive and constructs its arithmetic subalgebra for every
number field. These were unconditional. For our $K$, A makes
the nontrivial zeros of its partition function
$\zeta_K=\zeta L(\,\cdot\,,\chi_{-3})$ lie on the critical
line. It does not newly construct this dynamical system.

\begin{proposition}[AC-K07: a character observable is not a positive trace]
\label{AC-K07}
On $\ell^2(I_K)$, where $I_K$ is the set of nonzero integral
ideals, set $H e_{\mathfrak a}=\log N\mathfrak a\,e_{\mathfrak a}$.
For a finite-order Hecke character $\chi$ of conductor
$\mathfrak f$, extend $\chi(\mathfrak a)$ by zero when
$(\mathfrak a,\mathfrak f)\ne1$ and define
$D_\chi e_{\mathfrak a}=\chi(\mathfrak a)e_{\mathfrak a}$.
For $\Re s>1$,
\[
 \operatorname{Tr}(e^{-sH})=\zeta_K(s),\qquad
 \operatorname{Tr}(D_\chi e^{-sH})=L_K(s,\chi),\qquad
 |L_K(s,\chi)|\le\zeta_K(\Re s).
\]
For nontrivial cubic $\chi$, $D_\chi$ is not positive and
$\Re D_\chi$ has eigenvalue $-1/2$. Positivity of the
Bost--Connes Gibbs states therefore does not, by this trace
identity, supply the positivity of the cubic Weil form.
\end{proposition}
\begin{proof}
For $\Re s>1$ the sum of the absolute diagonal entries is
bounded by the absolutely convergent ideal zeta series. This
proves trace class and the identities by the definition of
the primitive Hecke series. The bound is the triangle
inequality. A faithful cubic ray-class character takes the
values $1,\omega,\omega^2$ on ideals prime to its conductor,
because every ray class is represented by such an ideal.
Its diagonal operator consequently has nonreal eigenvalues,
and its real part has eigenvalue $\Re\omega=-1/2$.
The assertion excludes this specific positivity transfer;
it is not a counterexample to cubic RH.
\end{proof}

\begin{proposition}[AC-K08: the cubic tensor bootstrap retains its target]
\label{AC-K08}
Let $V$ be the two-dimensional standard representation of
$S_3$ in AM-H01, and let $\varepsilon$ be its sign character.
For every integer $m\ge1$,
\[
 V^{\otimes m}\simeq
    1^{\oplus a_m}\oplus\varepsilon^{\oplus a_m}
                           \oplus V^{\oplus c_m},\qquad
 a_m=\frac{2^m+2(-1)^m}{6},\quad
 c_m=\frac{2^m-(-1)^m}{3}.
\]
In particular the coefficient of $V$ is positive for all $m$.
For even $m$, all Frobenius traces of this tensor power are
nonnegative, but its Artin $L$-function is still
\[
 L(s,V^{\otimes m})=\zeta_K(s)^{a_m}L_p(s)^{c_m}.
\]
Under A, its critical-line assertion is exactly equivalent
to that of $L_p$. Passing to an even tensor power therefore
does not eliminate the outstanding sector.
\end{proposition}
\begin{proof}
The three conjugacy classes have sizes $1,3,2$, and $V$ has
character values $2,0,-1$. Taking inner products of
$(2^m,0,(-1)^m)$ with the three irreducible characters gives
the stated multiplicities. They are nonnegative integers
(for $m=1$, $a_1=0$), and $c_m>0$. At a ramified prime the
same representation identity holds before taking inertia
invariants, so Artin multiplicativity gives the displayed
full $L$-function identity there too. All its factors have
no poles in the open critical strip, which proves the last
assertion. At a ramified prime, with inertia group $I$ and a
Frobenius lift $g$, the trace on invariants is
$|I|^{-1}\sum_{h\in I}\operatorname{tr}V(gh)^{m}$, which is
nonnegative for even $m$ as well. Nonnegative local traces of even powers are
therefore not the missing global one-sided window estimate.
\end{proof}

\paragraph{The commuting phase in geometric language.}
Over $U=\operatorname{Spec}\mathcal O_K[1/(3p)]$, the equation
$z^3=p$ is a finite \'etale $\mu_3$-torsor. Its two nontrivial
characters give rank-one lisse Artin sheaves, pure of local
weight zero. The action of $\mu_3$ commutes with Frobenius
over $K$, so the individual character projectors exist over
$\mathbb Q(\omega)$ or $\mathbb Q_\ell(\omega)$. Over
$\mathbb Q$ only the sum of the two nontrivial projectors is rational.
The integral failure of these projectors at three is already
computed in AM-H06--AM-H08; it is not repaired by a change
of vocabulary. For the CM Fermat cubic $x^3+y^3=1$, the
automorphism $(x,y)\mapsto(\omega x,y)$ instead acts on
$H^1$, giving weight-one Hecke realizations. Their local
weight theorem and their angular types are different data
from RH of their number-field Euler products. The relevant
frameworks are SGA~1, expos\'e~XI, \S6 (Kummer torsors),
and SGA~6, expos\'e~V, 2.2(b) and Lemma~2.2.1 (exterior
powers and representation rings). No inequality in these
framework statements implies H04.4 merely from the phase.

\subsection{An unconditional cubic window reached by the operator route}

\begin{proposition}[AC-K09: strict cubic positivity through $\lambda=4$]
\label{AC-K09}
For $p=269$, let $W_{\chi,\lambda}$ be the Weil quadratic form
of the primitive completed $L_K(s,\chi_{269,3})$ on logarithmic
test functions supported in $[-\log\lambda,\log\lambda]$.
With Fourier normalization as in AC-T05, for every
$1<\lambda\le4$ and every vector in its closed form domain,
\[
                     W_{\chi,\lambda}(v)\ge\|v\|_2^2.
\]
This assertion uses no RH for the cubic character and no input A.
It settles this finite range of windows, not all windows.
\end{proposition}
\begin{proof}
The primitive conductor in AC-K04 gives
$Q=|D_K|N\mathfrak f=217083$ and completed factor
$Q^{s/2}\Gamma_{\mathbb C}(s)L_p(s)$, where
$\Gamma_{\mathbb C}(s)=2(2\pi)^{-s}\Gamma(s)$.
The explicit formula, proved by the same logarithmic-derivative
contour calculation as AC-T05, has archimedean symbol
\[
 a(t)=\log Q-2\log(2\pi)+2\Re\psi(1/2+it).
\]
There is no pole term, since the character is nontrivial.
The convergent digamma difference series gives
$\Re\psi(1/2+it)\ge\psi(1/2)=-\gamma-2\log2$.
Thus $a(t)\ge c_0=\log Q-2\log(8\pi)-2\gamma$.

For the largest interval, $\lambda=4$, the only nonzero
prime-power coefficients with $n<16$ are
\[
\begin{array}{c|ccccc}
n&3&4&7&9&13\\ \hline
b_p(n)&\log3&2\log2&-\log7&\log3&-\log13.
\end{array}
\]
The prime above three splits by AC-K04. At inert rational
primes only even powers contribute, as in AM-H02. At seven
and thirteen, $269^{(\ell-1)/3}\not\equiv1\pmod\ell$, so
AM-H02 gives the two negative entries. These assertions
also follow by direct residue computations ($3^2\not\equiv1$
modulo seven and $9^4\not\equiv1$ modulo thirteen).
The translation by $\log16$ has zero compression at the
endpoint and contributes nothing.

For a shift $a$ in an interval of length $T$, AC-T14's
finite-chain decomposition gives
$\|T_a+T_{-a}\|=2\cos(\pi/(\lceil T/a\rceil+1))$
when $a<T$, and zero when $a\ge T$. At $T=\log16$, the
maximal chain lengths are respectively $3,2,2,2,2$ for
the five displayed entries. Consequently the norm of the
whole prime operator is at most
\[
 B_4=\sqrt2\,\frac{\log3}{\sqrt3}+\log2
       +\frac{\log7}{\sqrt7}+\frac{\log3}{3}
       +\frac{\log13}{\sqrt{13}}.
\]
The exact rational-interval certificate
\path{scripts/check_cubic_269.py} proves $c_0-B_4>1$.
Its elementary enclosures use the positive arctanh series
for logarithms with its geometric remainder, the alternating
Machin arctangent series for $\pi$, integer square-root
bounds, and $H_{100}-\log101<\gamma<H_{100}-\log100$.
All comparisons in that certificate are rational, not
floating-point sign decisions. The explicit formula therefore
gives the stated inequality first on smooth test functions.
It extends by closure to the full logarithmic-symbol form
domain. Restriction to a smaller support interval preserves
both the form and the norm, proving all $\lambda\le4$.
\end{proof}

\subsection{The exact gate retained from the newest checkpoint}

Write $b_p(n)$ for the full coefficients of $-L_p'/L_p$ and
$D_p=L_p'/L_p$. AM-H04 defines
\[
 G_p(y)=\sum_{n\le y}b_p(n)(y/n-1)+yD_p(1)
\]
and verifies every input of the generic theorem
\texttt{Deduction.landau\_window} in
\path{Development/AutomorphicArithmetic/DedekindLFunctions/LandauTheorems/landauWindow.lean}.
The recovered source of that theorem is present inside the
newest checkpoint's inherited source directory and was read
in this run; its dependencies were not rebuilt in Lean.
The remaining premise is, for each $\theta>1/2$, a constant
$C_\theta$ and arbitrarily large $Y$ such that
\begin{equation}
 G_p(y)\le C_\theta y^\theta
             \quad(Y\le y\le Y^{1+\delta}),\qquad\delta>0.
                                                               \label{eq:H044}
\end{equation}
None of AC-K01--AC-K09 proves \eqref{eq:H044}; the window in
AC-K09 is a bounded multiplicative test window and must not
be confused with an unbounded arithmetic ramp window in $y$.

The minimal-conductor host above removes ambiguity about the
finite algebra and its ramified Frobenius lifts. It supplies
no inequality on the individual twisted cohomology lengths
or the growing-host costs in ALP\_MACHINE's final Q3 gate.
The stronger full trace formula over $K$ would require the
Weil forms for \emph{all} its unitary Hecke characters to be
positive. AC-K05 proves the norm sectors. The cubic sector
and the nonzero angular sectors remain separate analytic
questions. Their failure to follow from the constructions
listed here is a missing proof, not a theorem of logical
independence from A.
